%\documentclass[sigconf]{acmart}
% 审稿版可改为以下两种之一（以具体 CFP 为准）：
%\documentclass[manuscript,screen,review]{acmart}
\documentclass[sigconf,review,anonymous]{acmart}

\AtBeginDocument{%
  \providecommand\BibTeX{{%
    Bib\TeX}}}

% =========================
% Packages
% =========================
\usepackage{booktabs}
\usepackage{multirow}
\usepackage{amsmath}
\usepackage{graphicx}
\usepackage{subcaption}

% =========================
% Rights / Conference info
% 投稿审稿阶段通常可先注释或按会议要求保留
% 录用后再替换成会议提供的正式信息
% =========================
% \setcopyright{acmlicensed}
% \copyrightyear{2026}
% \acmYear{2026}
% \acmDOI{XXXXXXX.XXXXXXX}
% \acmConference[Conference Acronym '26]{Full Conference Name}{Month DD--DD, 2026}{City, Country}
% \acmISBN{978-1-4503-XXXX-X/2026/XX}

% 如果投稿系统提供 submission ID，可打开
% \acmSubmissionID{123-A56-BU3}

% =========================
% Title
% =========================
\title[Constrained MORL for Cyber Defense]{Constrained Multi-Objective Reinforcement Learning for Autonomous Cyber Defense in CybORG}

% =========================
% Authors
% 审稿双盲时请改成 anonymous 版本，或按会议要求处理作者信息
% =========================
\author{First Author}
\affiliation{%
  \institution{Your University}
  \city{Your City}
  \country{Your Country}}
\email{first.author@xxx.edu}

\author{Second Author}
\affiliation{%
  \institution{Your University}
  \city{Your City}
  \country{Your Country}}
\email{second.author@xxx.edu}

\renewcommand{\shortauthors}{First Author et al.}

% =========================
% Abstract
% 这是一版“能直接用来继续改”的 starter abstract
% =========================
\begin{abstract}
Autonomous cyber defense requires sequential decision making under competing operational objectives. 
A defender must contain attacker progress while limiting business disruption and defense overhead, making the problem naturally multi-objective and constrained. 
In this paper, we study constrained multi-objective reinforcement learning for autonomous cyber defense in the CybORG environment. 
We formulate blue-team defense as a partially observable sequential decision problem with three objectives: security effect, business impact, and defense cost, while enforcing operational constraints on unacceptable side effects. 
We propose a two-stage learning framework that first discovers an initial set of preference-diverse defense policies and then expands the Pareto set under constraints. 
We evaluate the method against scalarized and constraint-handling baselines using Pareto-front and utility-oriented metrics, including hypervolume and expected utility, together with cyber-semantic metrics such as critical asset protection and attack-impact suppression. 
Results are expected to show that the proposed approach achieves a better trade-off frontier and stronger policy assignment quality across diverse defense preferences.
\end{abstract}

% =========================
% CCS Concepts
% 这里先放占位，正式投稿前请用 ACM CCS 工具生成并替换
% =========================
\begin{CCSXML}
<ccs2012>
 <concept>
  <concept_id>10002978.10002991.10002994</concept_id>
  <concept_desc>Security and privacy~Intrusion/anomaly detection and malware mitigation</concept_desc>
  <concept_significance>500</concept_significance>
 </concept>
 <concept>
  <concept_id>10010147.10010257.10010293.10010294</concept_id>
  <concept_desc>Computing methodologies~Reinforcement learning</concept_desc>
  <concept_significance>500</concept_significance>
 </concept>
</ccs2012>
\end{CCSXML}

\ccsdesc[500]{Security and privacy~Intrusion/anomaly detection and malware mitigation}
\ccsdesc[500]{Computing methodologies~Reinforcement learning}

% =========================
% Keywords
% =========================
\keywords{
autonomous cyber defense,
constrained multi-objective reinforcement learning,
Pareto optimization,
CybORG,
network security
}

\begin{document}

\maketitle

% =========================
% 1. Introduction
% =========================
\section{Introduction}
Autonomous cyber defense is inherently a multi-objective decision-making problem. 
A blue agent must slow or stop attacker progression, preserve business continuity, and limit the operational cost of defensive actions. 
Optimizing only one objective often leads to undesirable trade-offs, such as aggressive containment with excessive disruption or low-cost defense with insufficient protection.

Recent progress in reinforcement learning has made it possible to learn adaptive defense policies in cyber environments. 
However, most prior work either optimizes a single scalar reward or handles constraints without explicitly modeling the full Pareto trade-off among competing defense objectives. 
This gap is particularly important in cyber operations, where different organizations may prioritize security, business continuity, and intervention cost differently.

In this work, we investigate constrained multi-objective reinforcement learning (C-MORL) for autonomous cyber defense in CybORG. 
Our setting focuses on a blue defender interacting with an adversarial red process under partial observability. 
We model three objectives: security effect, business impact, and defense cost. 
On top of this multi-objective formulation, we impose operational constraints to avoid unacceptable defense behavior.

The main contributions of this paper are as follows:
\begin{itemize}
    \item We formulate autonomous cyber defense as a constrained multi-objective sequential decision problem with security, business, and cost objectives.
    \item We adapt a two-stage Pareto-oriented learning framework to the CybORG cyber-defense setting.
    \item We design a cyber-specific evaluation protocol combining Pareto metrics, utility metrics, and semantic defense metrics.
    \item We empirically compare the proposed method with scalarization and constraint-focused baselines.
\end{itemize}

% =========================
% 2. Background and Related Work
% =========================
\section{Background and Related Work}
\subsection{Autonomous Cyber Defense}
Briefly introduce RL-based cyber defense, CybORG, and why partial observability matters.

\subsection{Multi-Objective Reinforcement Learning}
Introduce MORL, Pareto front discovery, preference-conditioned methods, and policy assignment.

\subsection{Constrained Reinforcement Learning}
Introduce operational constraints, Lagrangian/CPO-style ideas, and why constraints matter in cyber defense.

\subsection{Positioning of This Work}
Explain that your paper focuses on transferring and adapting constrained MORL to a cyber-defense environment, rather than solving a fully general game-theoretic equilibrium problem.

% =========================
% 3. Problem Formulation
% =========================
\section{Problem Formulation}
We consider autonomous cyber defense as a partially observable decision problem. 
At each time step, the blue defender observes a partial view of host and network status and selects a defensive action from a discrete action set. 
The environment evolves under the interaction between defender actions and attacker behavior.

\subsection{Observation, Action, and Transition}
Define observation $o_t$, action $a_t$, latent state $s_t$, and transition dynamics.

\subsection{Multi-Objective Return}
We define a vector reward
\begin{equation}
\mathbf{r}_t = \bigl[r_t^{\text{sec}},\; r_t^{\text{biz}},\; r_t^{\text{cost}}\bigr],
\end{equation}
where $r_t^{\text{sec}}$ captures security effect, $r_t^{\text{biz}}$ measures business impact, and $r_t^{\text{cost}}$ measures defense overhead.

The corresponding discounted return is
\begin{equation}
\mathbf{G}_t = \sum_{k=0}^{T-t} \gamma^k \mathbf{r}_{t+k}.
\end{equation}

\subsection{Operational Constraints}
Let $c_t$ denote a constraint signal encoding unacceptable side effects, such as excessive disruption or unsafe intervention frequency. 
The policy must satisfy
\begin{equation}
\mathbb{E}_{\pi}\left[\sum_{t=0}^{T} \gamma^t c_t \right] \leq \tau,
\end{equation}
where $\tau$ is a predefined budget or threshold.

\subsection{Preference-Conditioned Utility}
For evaluation, a preference vector $\mathbf{w}$ induces utility
\begin{equation}
U(\mathbf{G}; \mathbf{w}) = \mathbf{w}^{\top}\mathbf{G},
\end{equation}
with $\mathbf{w} \in \Delta^{2}$ for the three-objective case.

% =========================
% 4. Method
% =========================
\section{Method}
\subsection{Overview}
Our method follows a two-stage learning pipeline. 
Stage~1 learns an initial set of preference-diverse defense policies. 
Stage~2 expands and refines the policy set under explicit operational constraints, aiming to improve Pareto-front coverage while maintaining safe behavior.

\subsection{Stage 1: Preference-Diverse Policy Initialization}
Describe how multiple policies are trained under different preference vectors to obtain an initial approximation of the Pareto set.

\subsection{Stage 2: Constrained Pareto Expansion}
Describe how candidate policies are expanded/refined under the constraint mechanism. 
Explain whether the constraint is enforced through a penalty, Lagrangian update, projection, or another optimization component.

\subsection{Policy Selection / Assignment}
At test time, given a user preference vector $\mathbf{w}$, select the policy from the learned policy set that maximizes utility:
\begin{equation}
\pi^{*}(\mathbf{w}) = \arg\max_{\pi \in \Pi} U(\mathbf{J}(\pi); \mathbf{w}),
\end{equation}
where $\mathbf{J}(\pi)$ denotes the estimated objective vector of policy $\pi$.

% =========================
% 5. Environment and Experimental Setup
% =========================
\section{Environment and Experimental Setup}
\subsection{Cyber Environment}
State which CybORG scenario / wrapper / action space is used. 
Clarify whether the red side is fixed, scripted, sampled from a family, or learned.

\subsection{Reward Design}
Specify the exact definitions of:
\begin{itemize}
    \item security effect,
    \item business impact,
    \item defense cost,
    \item constraint signal(s).
\end{itemize}

\subsection{Baselines}
Include the most relevant baselines, for example:
\begin{itemize}
    \item Weighted-Sum PPO,
    \item Preference-Conditioned PPO,
    \item PCN,
    \item Lagrangian-PPO / constrained single-objective baseline,
    \item w/o constraint,
    \item w/o Pareto expansion.
\end{itemize}

\subsection{Evaluation Metrics}
We report both general MORL metrics and cyber-semantic metrics.
General metrics include Hypervolume (HV), Expected Utility (EU), and additional Pareto-set quality indicators if needed.
Cyber-semantic metrics include critical asset protection rate, red impact success rate, average business impact, and average defense cost.

\subsection{Implementation Details}
Provide training budget, seeds, optimizer, learning rate, rollout length, and hardware.

% =========================
% 6. Results
% =========================
\section{Results}

\subsection{Main Pareto and Utility Results}
Present the main comparison table on Pareto-front quality and utility.

\subsection{Constraint Satisfaction Results}
Present the main comparison table on constraint handling and safety behavior.

\subsection{Cyber-Semantic Evaluation}
Interpret the results from the perspective of practical defense outcomes.

\subsection{Ablation Studies}
Include ablations such as:
\begin{itemize}
    \item without constrained expansion,
    \item without Pareto extension,
    \item different preference distributions,
    \item different attacker settings.
\end{itemize}

% =========================
% 7. Discussion
% =========================
\section{Discussion}
Discuss why the proposed method works in cyber defense, under which preference regimes it is strongest, and where the main limitations remain.

% =========================
% 8. Conclusion
% =========================
\section{Conclusion}
This paper studies constrained multi-objective reinforcement learning for autonomous cyber defense in CybORG. 
By jointly modeling security effect, business impact, and defense cost under operational constraints, the proposed framework aims to provide a more realistic and decision-relevant defense policy set than single-objective alternatives. 
Future work will extend the approach to broader attacker uncertainty, richer constraints, and stronger multi-agent interaction models.

% =========================
% Figure example
% 用时取消注释
% =========================
% \begin{figure}[t]
%   \centering
%   \includegraphics[width=\linewidth]{figs/overview.pdf}
%   \caption{Overview of the proposed two-stage constrained MORL framework for autonomous cyber defense.}
%   \Description{A pipeline diagram showing stage-1 preference-diverse policy initialization, stage-2 constrained Pareto expansion, and test-time policy assignment.}
%   \label{fig:overview}
% \end{figure}

% =========================
% Table example
% 用时取消注释
% =========================
% \begin{table}[t]
%   \caption{Main Pareto and utility comparison.}
%   \label{tab:main_results}
%   \centering
%   \begin{tabular}{lccc}
%     \toprule
%     Method & HV $\uparrow$ & EU $\uparrow$ & Constraint Violation $\downarrow$ \\
%     \midrule
%     Weighted-Sum PPO & -- & -- & -- \\
%     PCN & -- & -- & -- \\
%     Lagrangian-PPO & -- & -- & -- \\
%     Ours & -- & -- & -- \\
%     \bottomrule
%   \end{tabular}
% \end{table}

% =========================
% Acknowledgments
% =========================
\begin{acks}
This work was supported by XXX.
We thank XXX for helpful discussions and infrastructure support.
\end{acks}

% =========================
% Bibliography
% =========================
\bibliographystyle{ACM-Reference-Format}
\bibliography{refs}

% =========================
% Appendix
% =========================
\appendix

\section{Additional Experimental Details}
Provide reward coefficients, constraint thresholds, scenario details, and training hyperparameters.

\section{Additional Results}
Provide extra tables, sensitivity analyses, and extended attacker-setting results.

\end{document}